\documentclass[12pt]{article}

\usepackage{amsmath}
\usepackage{amssymb}
\usepackage{amsthm}
\usepackage{geometry}
\usepackage{setspace}

\geometry{letterpaper, margin=1in}
\setstretch{1.15}

\newtheorem{theorem}{Theorem}
\newtheorem{definition}{Definition}
\newtheorem{remark}{Remark}

\begin{document}

\title{\textbf{Algebraic Descent in Mixed-Dimensional Gaussian Traces: \\ A Period Ring Formulation}}
\author{}
\date{}
\maketitle

\begin{abstract}
\noindent We construct an analytic invariant, the \textit{Mixed Gaussian Trace} $\mathcal{Z}(p, q)$, defined as the tensor product of the continuous Gaussian measure on $\mathbb{R}^p$ and the discrete heat trace on the flat torus $\mathbb{T}^q$, evaluated at the modular fixed point. By framing the trace within the Kontsevich-Zagier ring of periods and invoking Chudnovsky's theorem on the algebraic independence of $\pi$ and $\Gamma(1/4)$, we show that $\mathcal{Z}(p, q)$ undergoes a formal algebraic descent if and only if $2p = 3q$. At the fundamental integer solution $(p=3, q=2)$, the canonical period $\pi$ precisely cancels, and the invariant drops into the purely elliptic period ring governed by the complex multiplication of the lemniscate curve.
\end{abstract}

\vspace{0.5cm}

\section*{1. The Mixed Spectral Trace}
We define a strict analytic invariant via spectral geometry. Let $\Delta_{\mathbb{T}^q}$ be the positive Laplace-Beltrami operator on the standard flat torus $\mathbb{T}^q = \mathbb{R}^q / \mathbb{Z}^q$. The eigenvalues of the Laplacian are $\lambda_n = 4\pi^2 \|n\|^2$ for $n \in \mathbb{Z}^q$. The trace of the heat kernel at spectral time $t > 0$ is the multidimensional Jacobi theta function:
\begin{equation}
    K_{\mathbb{T}^q}(t) = \operatorname{Tr}\left(e^{-t\Delta_{\mathbb{T}^q}}\right) = \sum_{n \in \mathbb{Z}^q} e^{-4\pi^2 \|n\|^2 t}
\end{equation}

We evaluate this trace at the specific spectral time $t = \frac{1}{4\pi}$, corresponding to the self-dual point $\tau = i$ under the modular transformation $\tau \mapsto -1/\tau$. This yields exactly $K_{\mathbb{T}^q}\left(\frac{1}{4\pi}\right) = \theta_3(e^{-\pi})^q$. Concurrently, we evaluate the standard $L^2$-normalized Gaussian measure over $\mathbb{R}^p$.

\begin{definition}[The Mixed Gaussian Trace]
Let $\mathcal{Z}(p, q)$ be the tensor product of the continuous Gaussian measure on $\mathbb{R}^p$ and the discrete heat trace on $\mathbb{T}^q$ at $t = \frac{1}{4\pi}$:
\begin{equation}
    \mathcal{Z}(p, q) = \left( \int_{\mathbb{R}^p} e^{-\|x\|^2/2} \, d^px \right) \times \operatorname{Tr}\left(e^{-\frac{1}{4\pi}\Delta_{\mathbb{T}^q}}\right)
\end{equation}
\end{definition}

\section*{2. The Chowla-Selberg Formula and Periods}
We contextualize the evaluation of this trace within the theory of periods. The self-dual torus evaluated at $\tau = i$ corresponds to the elliptic curve $E: y^2 = x^3 - x$, which has Complex Multiplication (CM) by the Gaussian integers $\mathbb{Z}[i]$.

By the Chowla-Selberg formula, the periods of CM elliptic curves are evaluated analytically in terms of the Euler Gamma function at rational arguments. Specifically, the trace at $\tau=i$ resolves exactly to:
\begin{equation}
    \theta_3(e^{-\pi}) = \frac{\Gamma(1/4)}{2^{1/2}\pi^{3/4}}
\end{equation}

Taking the product across arbitrary dimensions $(p,q)$ yields the exact closed form:
\begin{equation}
    \mathcal{Z}(p, q) = (2\pi)^{\frac{p}{2}} \cdot \left(\frac{\Gamma(1/4)}{2^{1/2}\pi^{3/4}}\right)^q = 2^{\frac{p-q}{2}} \cdot \pi^{\frac{2p-3q}{4}} \cdot \Gamma\left(\frac{1}{4}\right)^q \label{eq:trace}
\end{equation}

This expression highlights that $\mathcal{Z}(p,q)$ is fundamentally a mixed state of two distinct transcendental periods: the circular period $\pi$ and the lemniscatic period associated with $\Gamma(1/4)$.

\section*{3. Algebraic Descent via Chudnovsky's Theorem}
In 1984, G. V. Chudnovsky proved a profound theorem in transcendental number theory: $\pi$ and $\Gamma(1/4)$ are algebraically independent over $\mathbb{Q}$. Consequently, no non-trivial polynomial relation exists between them. This dictates that the fractional exponent of $\pi$ in Equation \ref{eq:trace} cannot be algebraically simplified, absorbed, or rationalized by any power of $\Gamma(1/4)$.

We define the \textit{transcendental $\pi$-weight} as the rational exponent $\frac{2p-3q}{4}$. Setting this weight identically to zero is strictly mathematically necessary for the trace to algebraically descend from a mixed transcendental state over $\mathbb{Q}(\pi, \Gamma(1/4))$ into a pure algebraic extension of the lemniscate period.

\begin{theorem}[Period Descent Condition]
The Mixed Gaussian Trace $\mathcal{Z}(p,q)$ is algebraically independent of the circular period $\pi$ if and only if the dimensional parameters satisfy the linear Diophantine equation:
\begin{equation}
    2p = 3q
\end{equation}
\end{theorem}

\section*{4. The (3,2) Elliptic Trace}
The minimal non-trivial positive integer combination satisfying this period descent is $p=3$ and $q=2$. When evaluated at this dimensional ratio, the transcendental $\pi$-weight vanishes entirely ($\pi^0 = 1$):
\begin{align}
    \mathcal{Z}(3, 2) &= 2^{\frac{3-2}{2}} \cdot \pi^0 \cdot \Gamma\left(\frac{1}{4}\right)^2 \nonumber \\
    \implies \quad &\boxed{ \mathcal{Z}(3, 2) = \sqrt{2} \cdot \Gamma\left(\frac{1}{4}\right)^2 \approx 18.5899\dots }
\end{align}

\begin{remark}
The initial heuristic identity $G^* = \frac{\Gamma(1/4)^2}{\pi\sqrt{2}}$ evaluated a 1-dimensional continuous measure against a 2-dimensional toroidal trace ($p=1, q=2$). Under this rigorous formal framework, the appearance of $\pi$ in the denominator of $G^*$ is revealed to be an uncancelled period weight mandated by Chudnovsky's theorem: because $2(1) \neq 3(2)$, the invariant fails the algebraic descent condition, and the period $\pi$ is mathematically forced to remain. By mapping $\mathcal{Z}(3,2)$, we isolate the exact arithmetic resonance where the Euclidean volume measure perfectly commensurates with the modular weight of the heat trace.
\end{remark}

\end{document}
